% !TEX root = ../thesis.tex
\graphicspath{{Chapter3/Chapter3Figs/}{Chapter3/Chapter3Figs/}}


 
% \chapter{Incorporate Textual Information in Predicting Volatility Surface}\label{ch: chapter: vol_surface}

\chapter{Volatility Surface Prediction with Textual Information}\label{ch: chapter: vol_surface}

\section{Introduction}\label{sec:ch3: introduction}  

Modeling the Implied Volatility Surface (IVS) has long been a central challenge in financial economics and engineering, as volatility plays a crucial role in option pricing, risk management, and derivative valuation. Traditional IVS models, grounded in the Efficient Market Hypothesis (EMH; \cite{fama1970}), rely primarily on historical and quantitative data.

Early approaches, such as the polynomial parameterization proposed by \cite{dumas1998implied} and spline-based interpolations (\cite{kahale2004arbitrage}; \cite{wang2004no}), sought to represent the volatility smile and term structure across maturities. Stochastic volatility frameworks, including the Heston Model (\cite{Heston:1993}) and the Stochastic Alpha Beta Rho (SABR) model (\cite{sabr2002}), are typically calibrated to the observed volatility structure of the market.

More recent advances in machine learning (ML) have enabled the use of non-parametric and deep learning architectures that can learn complex patterns directly from data. For example, \cite{ackerer2020deep} applied deep neural networks to estimate IVS non-parametrically, \cite{medvedev2022multistep} employed Long Short-Term Memory (LSTM) networks to capture temporal dependencies, and \cite{kelly2023deep} utilized Convolutional Neural Networks (CNNs) to model spatial patterns across moneyness and maturity dimensions.

Although these models provide rigorous mathematical formulations, their reliance on purely numerical inputs limits their ability to capture qualitative information—such as economic narratives, market sentiment, and policy expectations—that often exert material influence on asset prices and volatility dynamics.

In contrast, textual information has increasingly been recognized as a valuable source of insight into market expectations, investor psychology, and macroeconomic uncertainty. Corporate disclosures, analyst reports, central bank statements, and financial news releases contain forward-looking signals that complement traditional numerical data.

A growing body of research demonstrates that text-derived information significantly influences financial markets. \cite{manela2017news} introduced the News Volatility Index to link media sentiment to equity market volatility, while \cite{bybee2024business} found that business news can forecast macroeconomic indicators. Similarly, \cite{vergote2012} showed that central bank communications affect the implied volatility of rate options. Other studies confirm that financial sentiment extracted from textual data can predict returns and volatility (\cite{schumaker2009quantitative}; \cite{hagenau2012automated}; \cite{ke2019predicting}; \cite{li2020incorporating}).

Early text-mining methods—such as bag-of-words (BoW) and dictionary-based sentiment analysis—were limited by their inability to capture semantic or contextual relationships, often oversimplifying linguistic meaning. The emergence of Large Language Models (LLMs) marks a transformative breakthrough in this regard. LLMs encode linguistic information into high-dimensional embeddings that capture semantics, tone, and contextual dependencies. By converting text into numerical representations, these models bridge the gap between qualitative and quantitative information. Empirical studies (\cite{chen2021stock}; \cite{ji2021stock}; \cite{chen2022expected}) show that embeddings generated by Transformer-based architectures, such as BERT and GPT, can enhance predictive accuracy when combined with market variables. This integration of textual and numerical data enables richer modeling of how qualitative narratives and sentiment shifts influence asset prices and volatility.

Building on these developments, this study proposes a novel framework that integrates LLM-derived textual embeddings with traditional quantitative market indicators to model and forecast the Implied Volatility Surface. By leveraging the representational power of LLMs, we aim to capture how textual sentiment, news content, and qualitative signals shape market expectations of volatility. The proposed framework extends beyond conventional curve-fitting and time-series methods, offering a new perspective on volatility forecasting that combines econometric rigor with linguistic intelligence.

This chapter is organized as follows. Section~\ref{ch3:sec:contribution} outlines the research gap and contributions; Section~\ref{ch3:sec:lr} reviews the literature on IVS modeling and the integration of textual data in financial forecasting; Section~\ref{ch3:sec:model} presents the proposed modeling framework; and Section~\ref{ch3:sec:data} describes the dataset used in this study.

%%% end of intro

\section{Research Question and Contribution}\label{ch3:sec:contribution}

While prior literature highlights the predictive value of both quantitative and textual information, existing studies on IVS modeling have not systematically integrated qualitative data into volatility surface prediction. This omission leaves a critical gap, as textual signals often convey forward-looking information that markets rapidly incorporate into option pricing dynamics. Therefore, we propose the following question in this research:

\begin{itemize}
\item \textbf{RQ1: Predictive Accuracy.} To what extent does the inclusion of textual embeddings improve the accuracy of IVS prediction compared with models based solely on quantitative inputs?
\item \textbf{RQ2: Model Comparison.} How do LLM-based embeddings perform relative to traditional text-mining approaches (e.g., bag-of-words or sentiment scores) in capturing the informational content of news and announcements?
\item \textbf{RQ3: Intraday Dynamics.} Can the proposed framework effectively capture intraday changes in the IVS immediately after major policy announcements and news events?
\item \textbf{RQ4: Practical Implications.} What are the implications of enhanced IVS prediction for option pricing, hedging strategies, and risk management in high-frequency trading environments?
\end{itemize}


This study addresses these questions by proposing a novel framework that integrates traditional quantitative market data with textual embeddings derived from large language models (LLMs) to model the Implied Volatility Surface (IVS). By leveraging the representational power of LLMs and neural networks in processing high-dimensional data, we aim to capture the immediate impact of news and policy announcements on volatility expectations, thereby extending the scope of IVS modeling beyond purely numerical inputs.

Unlike existing machine learning approaches, which predominantly focus on daily or lower-frequency forecasts, our research also emphasizes intraday prediction of the IVS, examining volatility surface dynamics immediately after news arrivals. This finer temporal resolution enables us to better capture market reactions in high-frequency settings and delivers practical value for traders, option market participants, and risk managers.



Overall, this study contributes to the literature on implied volatility modeling, financial econometrics, and machine learning in four key ways:

\begin{itemize}
\item \textbf{Textual Information in IVS Modeling:} We introduce qualitative data, such as news releases and policy announcements, as systematic inputs into IVS modeling, addressing a critical omission in the existing literature.
\item \textbf{Hybrid Modeling Framework:} We propose a novel framework that combines traditional quantitative market variables with LLM-based textual embeddings, bridging the gap between econometric models and modern NLP techniques.
\item \textbf{Benchmarking Against Established Models:} We benchmark the predictive performance of LLM-augmented models against traditional parametric (e.g., Heston, SABR) and non-parametric (e.g., spline-based) approaches to evaluate their relative strengths.
\item \textbf{Intraday Prediction of IVS:} Unlike most existing ML-based studies that focus on daily forecasts, we provide accurate intraday-level predictions of the volatility surface, enabling timely option pricing adjustments, improved hedging strategies, and more effective management of market microstructure risks.
\end{itemize}

To sum up, these contributions extend the methodological toolkit for volatility surface forecasting and highlight the potential of combining quantitative and qualitative data sources in advancing financial market prediction.


%%% end of contribution and question



\section{Literature Review}\label{ch3:sec:lr}

The implied volatility surface (IVS), derived from market option prices, is a critical tool for traders, who monitor it closely to provide quotes, manage risk, and value their portfolios. Consequently, there is a growing body of literature devoted to modeling and forecasting the implied volatility surface.

% traditional models

Traditional approaches to modeling the surface can be broadly categorized into parametric and non-parametric models.

Parametric models aim to characterize the IVS using interpretable parameters that capture its structural properties. A seminal example is the polynomial parametrization proposed by Dumas et al. (\cite{dumas1998implied}), which employs polynomial functions to represent the skewness and curvature of the volatility smile across different maturities. However, this approach assumes constant volatility and is limited in its ability to accurately forecast the full surface dynamics. To overcome these limitations, stochastic volatility models have been developed. Notable examples include the Heston Model (\cite{Heston:1993}) and the Stochastic Alpha Beta Rho (SABR) Model (\cite{sabr2002}), both of which model the underlying asset price and its volatility as stochastic processes. These models effectively capture volatility skewness across both moneyness and maturity dimensions. In addition, econometric models such as the GARCH model (\cite{Bollerslev:1986}) have proven successful in capturing empirical features of volatility, including persistence, mean reversion, and stochastic behavior.

In contrast, non-parametric models focus on fitting the IVS directly using flexible mathematical techniques without relying on interpretable parameters. Common methods include linear interpolation and cubic splines (\cite{kahale2004arbitrage}; \cite{wang2004no}), which are designed to produce smooth surfaces consistent with observed option prices. While non-parametric models offer greater flexibility and accuracy in fitting, they lack explanatory power—their parameters serve purely as curve-fitting tools without economic or probabilistic interpretation.



% ML models


Furthermore, with the rapid advancement of machine learning (ML) techniques across various domains, researchers have increasingly applied these methods to model the Implied Volatility Surface (IVS). Compared with traditional econometric and parametric models, ML-based approaches offer a notable advantage in handling high-dimensional datasets and capturing complex nonlinear relationships.


Early contributions, such as \cite{malliaris1996using} and \cite{guidolin2003option}, explored the use of neural networks and support vector machines to improve IVS prediction. More recent research has benefited from significant improvements in computational power, enabling the use of more complex neural architectures to capture the nuanced dynamics of the volatility surface. For example, \cite{ackerer2020deep} employed deep neural networks for non-parametric IVS modeling, \cite{medvedev2022multistep} utilized long short-term memory (LSTM) networks to incorporate temporal dependencies, and \cite{kelly2023deep} adopted convolutional neural networks (CNNs) to extract spatial patterns across moneyness and maturity dimensions.

In addition, the application of Generative Adversarial Networks (GANs) has recently emerged as a promising approach in financial risk management. For example, \cite{ge2025gan} shows the GAN model is outperformaned than other parametric models. \cite{vuletic2024volgan} further demonstrated that GAN-based models can incorporate the arbitrage-free conditions when generating the implied volatility surface. Moreover, \cite{cont2025tail} showed that GANs achieve remarkable performance in managing tail risk.



However, previous models only focus on the quantitative data, such as option prices, underlying assets' prices. Qualitative information—such as news articles, policy announcements, and official statements—can exert a significant influence on implied volatility. For example, the study by \cite{vergote2012} demonstrates how European Central Bank (ECB) announcements affect the implied volatility of 3M EURIBOR rate options. Despite this, neither traditional econometric models nor most machine learning approaches incorporate textual data as a predictive input in modeling the IVS. 

% This omission represents a critical gap in the literature, given that market sentiment and forward-looking policy signals conveyed through text can substantially alter volatility expectations.


% incorporate text data

On the other hand, in recent years, there has been growing interest in incorporating textual information into predictive models, particularly in the domain of stock price forecasting. Several studies have demonstrated the significant impact of news on financial markets. For instance, \cite{manela2017news} developed the News Volatility Index by mapping sparse news text features to option-implied volatility measures, while \cite{bybee2024business} provided evidence that business news can forecast macroeconomic indicators. In addition, a substantial body of literature has explored the extraction of sentiment signals from financial news and their predictive effects on market movements. Notable examples include \cite{schumaker2009quantitative}, \cite{hagenau2012automated}, \cite{ke2019predicting}, and \cite{li2020incorporating}, which investigate how sentiment derived from news articles can serve as a critical explanatory variable. Earlier research, such as \cite{schumaker2009textual} and \cite{de2014evaluating}, employed text-mining techniques—most notably, the bag-of-words approach—to quantify unstructured textual data for use in econometric models. In this study, this strand of literature is used as a methodological baseline: we adapt the n-gram frequency representation to our Dow Jones news and bond-option data, but evaluate it under the same FiLM-WGAN downstream model as the LLM embeddings rather than directly comparing with prior papers' reported results.


More recently, the emergence of Large Language Models (LLMs) has significantly improved the integration of textual data into prediction frameworks. LLMs are capable of encoding complex semantic information into high-dimensional embedding vectors, enabling more nuanced representation of textual inputs. Studies such as \cite{chen2021stock}, \cite{ji2021stock}, and \cite{chen2022expected} utilize these embeddings to transform unstructured financial news into structured numerical inputs. In particular, embeddings generated by BERT-style encoders and GPT-style decoders have been shown to enhance forecasting performance when used alongside traditional quantitative indicators.



Overall, the literature reflects a growing consensus that embeddings derived from Large Language Models (LLMs) provide a powerful modality for enhancing financial forecasting, particularly when integrated with traditional market variables. As LLM architectures and pretraining methods continue to advance, future research is expected to increasingly extend this methodology to a broader range of financial markets, including interest rate instruments and foreign exchange markets.



\section{Model Details}\label{ch3:sec:model}

Inspired by \cite{vuletic2024volgan} who applied a Generative Adversarial Network (GAN, \cite{goodfellow2020generative}) in calibrating the current Implied Volatility Surface (IVS), we adopt the GAN framework with conditional information to predict the IVS. Compared with alternative architectures such as Convolutional Neural Networks (CNNs) and Long Short-Term Memory (LSTM) networks, GANs have demonstrated superior accuracy in generating complex distributions.

To improve training performance, we employ the Wasserstein GAN (WGAN, \cite{arjovsky2017wasserstein}) instead of the standard GAN. Unlike the latter, which minimizes the Kullback–Leibler (KL) divergence, the WGAN reformulates the loss function in terms of the Wasserstein distance. This modification stabilizes training, enhances convergence, and yields smoother gradients, making it particularly well-suited for modeling the intricate structure of the IVS.


Furthermore, to ensure that the generated volatility surfaces remain consistent with financial theory, we explicitly impose non-arbitrage constraints, specifically butterfly arbitrage and calendar arbitrage, during the generation process. Following the works of \cite{sidogi2022creating}, \cite{vuletic2024volgan}, and \cite{na2023computing}, we apply the Conditional GAN (CGAN, \cite{mirza2014conditional}) and adjust the loss function to ensure that the generated surfaces are arbitrage-free.


Overall, the integration of conditional information with no-arbitrage constraints ensures that the resulting IVS is both accurate and consistent with financial theory, thereby preserving its economic interpretability and practical relevance.

Finally, our model will take the following parameters as our input:
\begin{itemize}
    \item the Implied Volatility Surface at time $t$.
    \item the Text-embedding vectors at time $t$.
\end{itemize}

The output will be:
\begin{itemize}
    \item the Implied Volatility Surface at time $t+1$.
\end{itemize}


\subsection*{Text-Embeddings}
% TODO: add how to get text embeddings here


\subsection*{Non-arbitrage Condition}
In an arbitrage-free market, the Implied Volatility Surface (IVS) must satisfy non-arbitrage conditions for both strikes and maturities. For a fixed maturity $T$, the call price $C(K,T)$ must be \emph{decreasing} in strike $K$ and \emph{convex} in $K$, ensuring that vertical spreads are non-negative and butterfly spreads are non-negative (Eq. \eqref{eq:ch3:monotonicity}). Equivalently, the risk-neutral density must satisfy $\frac{\partial^2 C}{\partial K^2} \geq 0$ (Eq. \eqref{eq:ch3:convexity}). For a fixed strike $K$, the call price must be \emph{non-decreasing} in maturity $T$, which precludes calendar-spread arbitrage (Eq. \eqref{eq:ch3:calendar}). Expressed in terms of the volatility surface, the total implied variance $w(T,k) = \sigma_{\text{imp}}^2(T,k)\,T$ should be non-decreasing in $T$, and as a function of log-moneyness $k$, it must generate convex option prices, thereby avoiding butterfly arbitrage (Eq. \eqref{eq:ch3:totalvariance}). 

% strikes (fixed maturity T)
\begin{equation}
\frac{\partial C(K,T)}{\partial K} \leq 0 
\quad \text{(Call price decreasing in $K$)} 
\label{eq:ch3:monotonicity}
\end{equation}

\begin{equation}
\frac{\partial^2 C(K,T)}{\partial K^2} \geq 0 
\quad \text{(Convexity in $K$, non-negative density)} 
\label{eq:ch3:convexity}
\end{equation}

% maturities (fixed strike K)
\begin{equation}
\frac{\partial C(K,T)}{\partial T} \geq 0 
\quad \text{(Call price non-decreasing in $T$)} 
\label{eq:ch3:calendar}
\end{equation}

% total implied variance w(T,k)
\begin{equation}
    \begin{split}
        w(T,k) &= \sigma_{\text{imp}}^2(T,k)\,T \\
        \frac{\partial w(T,k)}{\partial T}&\geq 0 
        \quad \text{(No calendar-spread arbitrage)} 
    \end{split}
    \label{eq:ch3:totalvariance}
\end{equation}

% Finally, we denoted the 










% \subsection*{GPU programming for Implied Volatility Calculation}\label{sec:ch3:iv_calculation}

% Processing large datasets remains a significant challenge, particularly in computing implied volatility. Several methods have been proposed to enhance computational efficiency, including Newton–Raphson, Brent’s method, and Artificial Neural Networks (see \cite{liu2019pricing}). In this study, we adopt the Newton–Raphson method (Eq.~\eqref{eq:ch3:newton} and Eq.~\eqref{eq:ch3:mse}) and implement it using GPU (Graphics Processing Unit) programming to accelerate computation.

% Unlike the Central Processing Unit (CPU), which is designed for general-purpose sequential processing, the GPU is optimized for handling a high volume of parallel operations. This makes it particularly well-suited for computational tasks involving large-scale numerical processing. Figure~\ref{fig: Chapter1: hardware} illustrates the architectural differences between CPUs and GPUs, highlighting the GPU’s capacity for parallelism that underpins its efficiency in financial computing tasks.

% \begin{figure}[ht]
%     \centering
%         \includegraphics[width=\textwidth]{Chapter1/Chapter1Figs/hardware.png}
%         \caption{CPU and GPU Architecture}\label{fig: Chapter1: hardware}
% \end{figure}

% Although modern CPUs have multiple cores, the number is still relatively small compared to GPUs. For instance, the AMD EPYC 9754 CPU has 128 cores\footnote{Source: \url{https://www.amd.com/en/products/cpu/amd-epyc-9754}}, whereas the Nvidia A30 GPU has 3,804 cores\footnote{CUDA cores, source: \url{https://www.pny.com/nvidia-a30}}. Consequently, GPUs can process many more tasks simultaneously than the CPU.


% Since GPUs are designed exclusively for numerical computations, the raw data (see Tab. \ref{tab: ch1: raw data}) must first be processed into detailed option data (see Tab. \ref{tab: ch1: option data}). Subsequently,this data is transferred to the GPU for the computation of the risk-neutral density. Figure \ref{fig: Chapter1: gpu_working} illustrates how data is processed by the GPU and CPU. Compared with CPU computing which takes days to process the data, the GPU takes 4 hours to finish calculate the data for a month. 


% \begin{table}[H]
%     \scriptsize
%     \centering
%         \begin{tabular}{*6{c}}
%             \toprule
%             \#RIC& Date-Time& Bid Price& Bid Size& Ask Price& Ask Size \\
%             \hline
%             TY1205R3& 2013-01-28T15:01:39.356235000Z& 0.0625&  4&         &     \\
%             TY1205R3& 2013-01-28T15:01:39.356235000Z& 0.0625&   &  0.09375& 100 \\
%             TY1205R3& 2013-01-28T15:05:13.016758000Z& 0.0625&   & 0.078125&     \\
%             TY1205R3& 2013-01-28T15:06:51.578188000Z& 0.0625&   & 0.078125&     \\
%             TY125O3& 2013-01-15T14:28:03.383112000Z&        &   & 0.015625&     \\
%             $\cdots$&                       $\cdots$&$\cdots$&  $\cdots$ & $\cdots$& $\cdots$  \\
%             TYM3& 2013-01-31T23:40:36.997409000Z& 130.046875& 14& 130.0625&    4 \\
%             TYM3& 2013-01-31T23:40:36.997409000Z& 130.046875& 14& 130.078125& 23 \\
%             TYM3& 2013-01-31T23:40:36.997409000Z& 130.046875& 14& 130.078125& 21 \\
%             TYM3& 2013-01-31T23:40:36.997409000Z& 130.046875& 17& 130.078125& 21 \\
%             \bottomrule
%     \end{tabular}  
%     \caption{Raw Data}
%     \label{tab: ch1: raw data}
% \end{table}

% \begin{table}[H]
%     \scriptsize
%     \centering
%         \begin{tabular}{*7{c}}
%             \toprule
%             \#RIC& Date-Time& Bid Price& Bid Size& Ask Price& Ask Size& Contract Type \\
%             \hline
%             TY10025B3& 2023-01-20T13:20:28.367299648Z& 15.0& 1.0& 15.109375& 1.0& Option \\
%             TY10025B3& 2023-01-20T13:20:30.643362911Z& 14.984375& 1.0& 15.109375& 1.0& Option \\
%             TY10025B3& 2023-01-20T13:20:33.058488285Z& 15.0& 1.0& 15.109375& 1.0& Option \\
%             TY10025B3& 2023-01-20T13:20:33.226949783Z& 14.984375& 1.0& 15.109375& 1.0& Option \\
%             TY10025B3& 2023-01-20T13:20:33.302647229Z& 15.0& 1.0& 15.109375& 1.0& Option \\
%             TY10025B3& 2023-01-20T13:20:44.047374299Z& 15.0& 1.0& 15.125& 1.0& Option \\
%             TY10025B3& 2023-01-20T13:20:45.423321752Z& 15.0& 1.0& 15.109375& 1.0& Option \\
%             TY10025B3& 2023-01-20T13:20:55.515039558Z& 15.0& 1.0& 15.125& 1.0& Option \\
%             TY10025B3& 2023-01-20T13:20:56.595216527Z& 15.015625& 1.0& 15.125& 1.0& Option \\
%             \bottomrule
%     \end{tabular}  
%     \caption{Option Data}
%     \label{tab: ch1: option data}
% \end{table}

% % \begin{figure}[ht]
% %     \centering
% %         \includegraphics[width=\textwidth]{Chapter1/Chapter1Figs/gpu_working.png}
% %         \caption{GPU Working}\label{fig: Chapter1: gpu_working}
% % \end{figure}



\subsection*{Conditional Generative Adversarial Network}

To predict the Implied Volatility Surface (IVS), we applied a Conditional Generative Adversarial Network (CGAN) with the Wasserstein-1 distance in the loss function during training. The network architecture is described through the generator-discriminator setup discussed below.

\subsubsection*{Generative Adversarial Network}
The Generative Adversarial Network (GAN) consists of two main parts, the Generator ($G(\cdot; \theta_g)$)and the Discriminator ($D(\cdot; \theta_d)$). The Generator will predict the IVS ($\sigma^G_{t+1}(m, T)$) based on the inputs, while the Discriminator will make a judgement of the predicted IVS ($\sigma^G_{t+1}(m, T)$) with the real IVS ($\sigma_{t+1}(m, T)$). $1$ will be given if the predicted IVS is matched with the real IVS, otherwise, the Discriminator will return $0$. A well-trained Generator can generate a similar IVS to the real one. 

\subsection*{Convolutional Neural Network (CNN) for Feature Extraction}
% TODO: add cnn here

\subsubsection*{Feature-wise Linear Modulation (FiLM) Layer for Conditional Information Encoding}
%  TODO: add FiLM here


\subsection*{FiLM-GAN Model Architecture}
% TODO: add details structure here



% where the $G(\cdot; \theta_g)$ is the Generator parameterized by $\theta_d$, which can be modelled using CNNs. the $D(\cdot; \theta_d)$ is the Discriminator parameterized by $\theta_d$. Wasserstein-1 distance is applied in calculating the loss for each iteration. 

% \begin{equation}\label{eq:ch3: wasserstein distance}
%     W_1(p_z, p_x) = \underset{\gamma \in \Gamma(p_z, p_x)}{inf} \, (\int \| x - y\| d\Gamma(x, y))
% \end{equation}






%

\section{The Dataset}\label{ch3:sec:data}  

In this section, we describe the datasets utilized in our research. The dataset comprises two primary components: (1) numerical data, including historical futures and options prices, and (2) textual data, consisting of relevant news articles. Our analysis focuses on the period from 17th March, 2022 to 27th July 2023, a time characterized by heightened economic uncertainty. During this interval, the Federal Reserve announced ten consecutive increases in the Federal Funds Rate (FFR), raising it from 0.50\% to 5.50\%. As the FFR serves as the benchmark rate for the U.S. economy, the Federal Reserve employs open market operations, such as the purchase and sale of government bonds, to guide the effective rate toward its target. These policy changes had a significant impact on the interest rate market. Accordingly, our research concentrates on this volatile period to investigate the effects of government announcements on the government bond option market dynamics. 




\subsection*{Futures and Options Dataset}\label{sec: ch3: dataset: option}
We use the U.S. 10-Year Treasury Note (T-Note) future contracts and the future option contracts. Both the future and the option contracts are listed in Chicago Mercantile Exchange (CME). All the contracts can be traded continuously except a one-hour halt from 4PM to 5Pm Central Time (CT). Hence, the contract can be regarded as trading for all the time allowing traders to reactions to economic events at almost anytime.

For the future contract, each contract has a face value of \$100,000 with a maturity of one-year. The contract is rolling, if one contract is expired, a new one-year contract will be opened automatically. There are 4 quarterly contracts listed with maturity months: March, June, September, and December. The contract is expired at the 7th business day before the end of maturity month. The contract is quoted on points with a basis of 100. 

The contracts are represented by 4 characters. For instance, ``TYM4\footnote{Refinitiv and Bloomberg use ``TY'' as the identifier for the U.S. 10-Year Treasury Note future contract, CME use ``ZB'' for the contract.}'' is the listed contract, where ``TY'' indicates that the underlying asset is the U.S. 10-year Treasury note, ``M'' denotes the expiration month (June),  and ``4'' signifies the last digit of the expiration year (2024).


The option contract is derived from one future contract but has maturities for each month. It is expired at the last business day of the maturity month. The option contract is American style which means traders can exercise the contract anytime before the maturity. Each option contract is identified by 6 to 8 characters. 

For instance, ``TY11225T3'' is the listed contract name. ``TY'' indicates that the underlying asset is the 10-year T-note future contract. ``11225'' represents the strike price of the contract, which is 112.25 points. ``T'' serves as the indicator of the contract maturity. Also it distinguishes between call and put options, different types of options have different sets of indicators for the contract maturity month. The example contract's maturity month is August. The final digital, ``3'', means the contract's maturity year is 2023. 

We extracted the limit order book data from the Refinitiv DataScope, a database vendor received data directly from the CME. Compared with the market order which is execute directly, although the order may not be executed, the limit order allows traders to submit specific prices about the asset based on their understanding about the market. Therefore, it can provide insights about their expectations about the asset. 

As mentioned previously, the US 10-Year T-Note future contract is the most liquid government bond future contracts in the market. The data size is extremely large. But we utilize the accessability of Refinitiv DataScope and extracted the future and option contract data by developing a Python SDK. Chapter \ref{ch:chapter:2} offers the engineering details about how we extract and process the data. 


Finally, we extracted the tick data for both the futures and options contracts covering the period from January 1, 2022, to December 31, 2023. The entire dataset occupies approximately 500 GB when stored in compressed files. Although data prior to January 1, 2016, is available in the database, it has not been extracted for this study. The current dataset suffices for our research objectives, additional data will be downloaded as necessary.



We utilize a 5-minute time interval to construct the Implied Volatility Surface (IVS). Since the option contracts are quoted in terms of price, we employ the Black Model (\cite{black76}, Eq.\ref{eq:ch3:black76}), which is the extended model for Black–Scholes model (\cite{BlackScholes:1973}) when the underlying asset is future contract ($F_t$), to invert the prices and compute the corresponding implied volatilities (IVs).


\begin{equation}\label{eq:ch3:black76}
    \begin{split}
        C_t^{bl}(F_t, r, T, \sigma) &= F_t N(d_1) - K e^{-r(T-t)} N(d_2) \\
        P_t^{bl}(F_t, r, T, \sigma) &= K e^{-r(T-t)} N(-d_2) - F_t N(-d_1)
    \end{split}
\end{equation}

\begin{equation*}
    \begin{split}
        d_1 &= \frac{\ln\left( \frac{F_t}{K} \right) + \frac{1}{2}\sigma^2(T-t)}{\sigma \sqrt{(T-t)}} \\
        d_2 &= d_1 - \sigma \sqrt{(T-t)}
    \end{split}
\end{equation*}


Where, $C_t^{bl}(\cdot)$ and $P_t^{bl}(\cdot)$ are the theoretical option prices, $T$ is the maturity, $r$ is the risk-free rate, $\sigma$ is the corresponding implied volatility (IV), $N(\cdot)$ is the cumulative density function (CDF) for standard normal distribution. Therefore, for a given market price $C_{i,t}^{mkt}$ or $P_{i,t}^{mkt}$, we can extract the IV ($\sigma_{i,t}^{*}$) by minimized the Mean Squared Error between the market price ($C_{i,t}^{mkt}$) and the theoretical price ($C_{i,t}^{bl}$):

\begin{equation}\label{eq:ch3:mse}
    \begin{split}
        \sigma_{i,t}^{*} &= \arg\min_{\sigma} \left( C_{i,t}^{mkt} - C_{i,t}^{bl}(F_t, r, T, \sigma) \right)^2  \\
        s.t. &\quad \sigma_{i,t}^{*}>0
    \end{split}
\end{equation}


Given the monotonicity and differentiability of the Black Formula with respect to volatility, we apply the Newton–Raphson method (Eq.\ref{eq:ch3:newton}) to efficiently solve for the implied volatility through numerical optimization.

\begin{equation}\label{eq:ch3:newton}
    \sigma_{n+1} = \sigma_n - \frac{f(\sigma_n)}{f'(\sigma_n)} = \sigma_n - \frac{C^{bl}_{i,t}(F_t, K, r, T, \sigma_n) - C^{mkt}_{i,t}}{Vega(\sigma_n)}
\end{equation}

Where $f(\cdot)$ is the objective function (Eq.\ref{eq:ch3:mse}), $Vega(\cdot)$ is the \textit{Vega} for this option contract, which is given by:

\begin{equation*}
    \begin{split}
        Vega(\sigma) 
        &= \frac{\partial C}{\partial \sigma} \\
        &=  e^{-r(T-t)} F_t \sqrt{(T-t)} \cdot \phi(d_1)
    \end{split} 
\end{equation*}

Here, $\phi(\cdot)$ denotes the standard normal probability density function (PDF). The $Vega$ calculation is the same for both the call option and the put option. 

\begin{figure}
    \centering
    \includegraphics[width=\textwidth]{Chapter3/Chapter3Figs/vol_example/future_2022-09-21.png}
    \caption{\textbf{Futures Price Movements around FOMC Press Release}, \textbf{TYH3} is the futures contract derived from the U.S. 10-Year Treasury Note with maturity 2033-03-15. Significant jumps can be observed around the announcement.}
    \label{fig:ch3:future_example}
\end{figure}



\begin{figure}[htbp]
\centering
\begin{minipage}[t]{0.47\linewidth}
\includegraphics[width=\textwidth]{Chapter3/Chapter3Figs/vol_example/vol_2022-10-31.png}
\end{minipage}
\hfill
\begin{minipage}[t]{0.47\linewidth}
\includegraphics[width=\textwidth]{Chapter3/Chapter3Figs/vol_example/vol_2022-11-30.png}
\end{minipage}
\caption{\textbf{Implied Volatility Surface before and after FOMC announcement}. The Implied Volatility is calculated by the average ask-bid price of Out-of-The-Money (OTM) call options and put options in 5-minute intervals.}
\label{fig:ch3:vol_example}
\end{figure}



Finally, for example, at 2:00 PM Eastern Daylight Time (EDT) on November 2, 2022, the FOMC announced an increase in the Federal Funds Rate (FFR) by 75 basis points, raising from 3.00\% to 3.75\% (lower bound). Figure~\ref{fig:ch3:future_example} illustrates the subsequent movements in futures prices, while Figure~\ref{fig:ch3:vol_example} displays the Implied Volatility Surfaces (IVS) for U.S. 10-Year Treasury Note option contracts with maturities on October 31, 2022, and November 30, 2022, respectively, following the announcement. The rate hike had a substantial impact on both the futures and options markets, reflecting heightened volatility and adjusted market expectations in response to the policy shift.




\subsection*{Textual Dataset}\label{sec:ch3:dataset:text}

% We use the daily number of headlines reported by Dow Jones as our primary
% measure of public information.' Dow Jones is the largest U.S. business
% newswire source, maintaining five wire services in addition to publishing the
% Wall Street Journal and Barron's. As described in Sommer (1984), Thomp-
% son, Olsen, and Dietrich (1987), and Patell and Wolfson (1982), most of the
% information compiled by Dow Jones originates from publicly traded compa-
% nies, as the major stock exchanges require member firms to provide all
% material information to Dow Jones in a timely fashion. The bulk of the
% information processed by Dow Jones is made public via the Wall Street
% Journal and the Dow Jones News Service, the latter commonly referred to as
% the Broadtape. For the period we analyze, Dow Jones operated the Broadtape
% between the hours of 7:30 A.M. and 7:00 P.M. (EST) each business day


For the textual dataset, we utilize news reports from ``Dow Jones Newswires'' as our primary source. Operated by Dow Jones \& Company, Dow Jones Newswires is one of the largest newswire services in the United States. It provides real-time information on market movements, government announcements, corporate disclosures, and other key economic developments, covering major segments such as equity markets, interest rate markets, and foreign exchange markets. Notably, the service offers millisecond-level latency, making it a valuable textual resource for financial research, as demonstrated in prior studies such as \cite{thompson1987attributes}, \cite{mitchell1994impact}, \cite{tetlock2008more}, and \cite{engelberg2012shorts}.

Accordingly, for our study, we extract articles specifically related to the ``Federal Reserve'' from this data source. The news is archived in a structured format, where each entry includes a headline, publication date and time, and a content summary. An illustrative example is shown in Table~\ref{tab:ch3:example_news}.


\begin{table}[htbp]
    \centering
    \footnotesize
    \begin{tabular}{l p{11cm}}
        \toprule
        \textbf{Indicator}& \textbf{Texts} \\
        \midrule
        HD& 	The Fed, State Taxes, Tesla, Roundtable Picks: Community Conversations -- Barrons.com \\
        PD& 	2 February 2023 \\
        ET& 	22:04 \\
        SN& 	Dow Jones Institutional News \\
        LP&     The Fed continues to tap the brakes but is doing so with a lighter touch as it seeks to tame inflation without pushing the economy into a ditch. 

        As expected, the Federal Open Market Committee raised interest rates Wednesday by 0.25 percentage points, bringing the target for the federal-funds rate to a range of 4.5\% to 4.75\%.\\
        \bottomrule
    \end{tabular}
    \caption{Example News from Dow Jones Newswires}
    \label{tab:ch3:example_news}
\end{table}



In this format, ``HD'' denotes the headline of the news article, ``LP'' refers to the lead paragraph or summary, while ``PD'' and ``ET'' indicate the publication date and exact time of release, respectively. Unlike approaches that use the full text of the article, we utilize only the summary (``LP'') in our model, as it typically contains the most salient and informative content relevant to financial market reactions. Our final dataset contains 14,872 relevant news from 27th Jan 2022 to 30th Dec 2023. 



% count    14900.000000
% mean       118.868658
% std         62.518772
% min          2.000000
% 25%         61.000000
% 50%        126.000000
% 75%        171.000000
% max        387.000000

% count    14900.000000
% mean        16.215839
% std          4.618225
% min          4.000000
% 25%         12.000000
% 50%         16.000000
% 75%         19.000000
% max         46.000000
% Name: hd_token_count, dtype: float64



% After that, we employed a ``tokenizer'' converting the textual news into vectors of tokens. The ``tokenizer'' is a pre-trained dictionary mapping a ``word'' with a specific ``number''. Noticeably, to maximized the efficiency and reduce the size of the ``dictionary'', the ``word'' sometimes refers to a part of a word, for example, ``negative'' may be split into ``nega'' and ``tive''. Moreover, the ``tokenizer'' encode not only the words but also the content. 

Next, we applied a ``tokenizer'' to convert the news text into sequences of tokens. The tokenizer is a pre-trained dictionary that maps ``word'' into a integer identifier. To maximize efficiency and reduce dictionary size, most modern tokenizers operate at the ``sub-word level'', e.g., Byte Pair Encoding or WordPiece. For instance, a word like “negative” may be split into "nega" and "tive". Table~\ref{tab:ch3:summary_news} summarized the number of tokens for the ``Headline (HD)'' and the ``Line Paragraph (LP)''.
% 2022-01-27 00:00:00 to 2023-12-30 07:38:00

% TODO: add summary stats here, 20250625

\begin{table}[htbp]
\centering
\footnotesize
\begin{tabular}{cccccccc}
\toprule
Name& Mean & Std & Min & 25\% & 50\% & 75\% & Max \\
\midrule
HD& 16.2158 & 4.6182 & 4 & 12 & 16 & 19 & 46 \\
LP& 118.8687 & 62.5188 & 2 & 61 & 126 & 171 & 387 \\
\bottomrule
\end{tabular}
\caption{Summary Statistics for HD and LP.}
\label{tab:ch3:summary_news}
\end{table}
% TODO: 20250625, add Bert here.

% On the other hand, to incorporate the textual information into our model, inspired by \cite{7550882}, \cite{ding2015deep}, we employed a BERT model, encoding the texts into a vector of tokens. The BERT model can not only encode the position information but also the content information into the token. Therefore, the vector will provide a comprehensive repersentative of the textual information. 

% Furthermore, to incorporate textual information into our model, and following \cite{7550882} and \cite{ding2015deep}, we employ an embedding model from Large Language Model (LLM) to transform each document into contextual representations. The model will convert the token vector into an embedding vector. The embedding vector is a vector of numbers encoded both semantic content and positional information, which provided an accurate and quantitative measurement of the textual data. 

% An embedding is a vector (list) of floating point numbers.

Furthermore, to integrate textual information into our model, inspired by \cite{7550882} and \cite{ding2015deep}, we employ an LLM-based embedding model to map each news article into contextual representations. In this process, the model transforms the token sequence into a fixed-length embedding vector which captures both semantic content and word order through positional encodings. Since LLM-based embedding models are pre-trained on large scale corpora, they can provide comprehensive and informative representations of text, making them well-suited for downstream financial analysis.


In this research, we employ the ``text-embedding-3-large'' model developed by OpenAI (\cite{test_embedding2021}). We experiment with both 512-dimensional and 1,024-dimensional embedding vectors to evaluate their effectiveness in representing textual information. The 1,024-dimensional embeddings offer richer representational capacity, potentially capturing more nuanced semantic and contextual features of the text, while the 512-dimensional embeddings provide a more computationally efficient alternative. By testing both, we assess the trade-off between representational richness and computational efficiency in our downstream tasks.




% We then aggregate these token embeddings into a fixed-length document vector (e.g., using the \texttt{[CLS]} representation or mean pooling). The resulting embedding provides a compact, quantitative representation of the text suitable for downstream prediction and econometric analysis.

% Tne advantage of using LLMs for text embedding is that they are extensively pre-trained on web-scale data already, which does not need the contrastive pre-training step used in existing state of the art text embedding models. 

% Concretely, BERT produces a sequence of context-dependent token embeddings by combining learned token, segment, and positional embeddings, thereby encoding both semantic content and word order. In our implementation, we use the pooled representation of the \texttt{[CLS]} token as a fixed-length document vector, which serves as a comprehensive representation of the text for downstream prediction.








% \subsubsection*{Information Adjustment}
% Since the option price contains forward-looking information while the future price is backward-looking, an adjustment need to be conducted to the future return distribution. We followed the idea proposed by \cite{barone2020option}. They create a link between the option-recovered physical density and the future-based realized density by:

% \begin{equation}\label{eq: ch3: barone}
%     \tilde{p}(S_t | \mathcal{F}_t) = w_t p^*(S_t| \mathcal{F}_t) + (1 - w_t)p(S_t | \mathcal{S}_t)
% \end{equation}

% where $\tilde{p}(S_t | \mathcal{F}_t)$ is the information-adjusted realized physical density, $p^*(S_t| \mathcal{F}_t)$ is the recovered physical density from the option risk-neutral density. $p(S_t | \mathcal{S}_t)$ is the realized physical density calculated from the future return, $w_t$ is the proportion of listed Deep Out of The Money options (DOTM) over the total number of listed Out of The Money options (OTM) and the DOTM options. Additionally, they suggested $w_t$ can be modeled by a function of the VIX index level. A large literature has proofed that the OTM options carry more information that the underlying assets, such as \cite{wei2010trading} and \cite{chakravarty2004informed}. Therefore, we take the idea by adding the weights to contract our convertor.







\section{Empirical Result}\label{sec: ch3: result}

\subsection*{Main Results}

\subsection*{Robustness Tests}



\section{Conclusion}\label{sec: ch3: conclusion}











